\subsection{Przykład rozmieszczenia w terminalu}
\fig{terminal.pdf}{1}{Rzut funkcjonalny terminalu 80 $\times$ 40 m. Kamery C01--C08, przejścia D01--D06, SOC i przykładowe trasy lokalne. Małe czerwone punkty pokazują 48 orientacyjnych miejsc detekcji SSP.}
Rzut przedstawia orientacyjne funkcje pomieszczeń, a nie rozstrzygnięty projekt architektoniczny. Sześć przejść kontrolowanych łączy strefę publiczną, pracowniczą, operacyjną i techniczną. Na wyjściach przejście od strony ewakuacji ma być możliwe bez karty, serwera i sieci.

Przyjęty przedmiar miedzi kat. 6A to \textbf{1800 m} łącznie z zapasami dla kamer i urządzeń IP. Przedmiar przewodów KD/SSWiN to \textbf{2200 m}, dobieranych oddzielnie do zasilania, RS-485 i linii nadzorowanych. Nie jest to jedna uniwersalna żyła dla wszystkich funkcji.

Dla pętli SSP: 600+300+150+100 = 1150 m; po dodaniu 20\% i zaokrągleniu przyjęto \textbf{1400 m}. Osobno przewidziano \textbf{400 m} na sygnalizację i sterowania, z podziałem obwodów według dobranego systemu. Typy kabli, wymagany czas zachowania funkcji i system mocowań ustala projekt ppoż.; nie wolno sprowadzać doboru do samego koloru powłoki.

\begin{note}{Następny poziom szczegółowości}
Po otrzymaniu rzeczywistych rzutów należy ustalić pomieszczenia, przeszkody, wysokości sufitów, drogi ewakuacji i wentylację. Następnie przeliczyć czujki, kadry, długości każdego obwodu, spadki napięć oraz macierz sterowań. Schematy w tym opracowaniu można edytować i wygenerować ponownie ze źródeł repozytorium.
\end{note}
